\section{Expressions and Processes}
\label{sec:processes}

This section introduces the notions of expressions and processes, algorithmic
typing for expressions, and operational semantics for expressions and processes
based on labelled transition relations.
%
% \paragraph*{Syntax}
%
The syntax of constants $\termc c$, values $\termc v$, expressions
$\termc e$ and processes $\termc p$ is defined by the
grammar below.
%
\begin{align*}
  \termc c \grmeq & \unitk
       % \grmor \reck
       \grmor \forkk 
       \grmor \newk
       \grmor \receivek 
       \grmor \sendk 
       \grmor \selecte C
       \grmor \waitk
       \grmor \terminatek
  \\
  \termc v \grmeq & \termc{c}
      \grmor \termc{x}
      \grmor \abse{x}{T}{e}
      \grmor \rece xTv
      \grmor \tabse\alpha\kind v
      \grmor \paire{v}{v}
      \grmor \rece xTv
      \grmor \tappe{\receivek}{T}
  \\& \tappe{\tappe{\receivek}{T}}{T}
      \grmor \tappe{\sendk}{T}
      \grmor \tappe{\tappe{\sendk}{T}}{T}
      \grmor \appe{\tappe{\tappe{\sendk}{T}}{T}}{v}
      \grmor \tappe{\selecte C}{\overline T}
  \\
  \termc e \grmeq & \termc{v}
       \grmor \appe{e}{e}
       \grmor \tappe{e}{T}
       \grmor \letu ee
       \grmor \paire ee
       \grmor \lete xxee
      \grmor
  \\& \matchwithe eCxeiI
  \\
  \termc p \grmeq & \expp e 
  \grmor \parp{p}{p}
  \grmor \newp{x}{y}{p} 
\end{align*}

All constructors are standard either in linear functional languages or in
session type languages. From functional languages we find function introduction
and elimination ($\abse xTe$ and $\appe{e_1}{e_2}$), linear pair introduction
and elimination ($\paire{e_1}{e_2}$ and $\lete xy{e_1}{e_2}$), linear unit
introduction and elimination ($\unitk$ and $\letu{e_1}{e_2}$). From session
languages we find most of the constructors in the form of constants, including
$\forkk$ to spawn a new thread, $\newk$ to create a new channel, $\receivek$ to
receive a value on a given channel, $\sendk$ to send a value on a given channel,
$\waitk$ to wait for a channel to be closed, $\terminatek$ to force channel
closing, and $\appe \selectk {C}$ to exercise a choice on a given
channel indicated by selector name $\constr$. The
operator to offer a choice on some channel cannot be captured by a type in our
type language, hence we have made it an expression ($\matchwithe eCxeiI$).
%
At the level of processes, we have expressions as threads $\expp e$,
parallel composition $\termc{p_1 \PAR p_2}$, and 
introduction of a channel by declaring its
two ends $\termc x, \termc y$: $\newp xyp$.

% \paragraph{Typing}

%\begin{figure}
\declrel{Context formation}{$\isCtx{\Gamma}{\kind}$}
\begin{mathpar}
  \infrule{Ctx-Empty}
  {}
  {\isCtx{\Empty}{\kind}}

  \infrule{Ctx-Ext}
  {\isCtx{\Gamma}{\kind}\\ \typeAgainst{T}{\kind} \\ \typec T = \Normal[+]{T}}
  {\isCtx{\Gamma, \ebind{x}{T}}{\kind}}
\end{mathpar}
\caption{Context formation}
\label{fig:ctx}
\end{figure}
\begin{figure}[t!]
  \declrel{Types for constants}{$\typeof c = \typec T$}
  \begin{align*}
    % \typeof \reck &= \TForall{\alpha}{\kindt}{\TForall{\beta}{\kindt}{((\alpha\rightarrow\beta)
    %                 \rightarrow (\alpha\rightarrow\beta)) \rightarrow (\alpha\rightarrow\beta)}}
    % \\
    \typeof \unitk & = \TUnit
    &
    \typeof \forkk & = \TFun[] {(\TFun[] \TUnit \TUnit)} \TUnit
    \\
    \typeof \newk & = \TForall\alpha\kinds{\TPair\alpha{\TDual \alpha}}
    &
    \typeof \closek &= \TFun[]{\TEnd}{\TUnit}
    \\
    \typeof \receivek &= \TForall\alpha\kindt{\TForall\beta\kinds
                        {?\alpha. \beta \to \TPair\alpha\beta}}
    \\
    \typeof \sendk &=
      \TForall\alpha\kindt{\TForall\beta\kinds
        {\TFun[]{\TPair{\alpha}{\TOut{\alpha}{\beta}}}{\beta}}}
  \end{align*}
  $
    \typeof {\appe \selectk {C_k}} = \TForall{\overline{\alpha}}{\overline\kindp}
                                     {
                                     {\TForall \beta \kinds {\TFun[]
                                     {\TOut{(\TPAppSet \proto {\{k\}} {\overline\alpha})}\beta}
                                     {\TFOut{\overline{T_k}}\beta}
                                     }}}
                                    \text{ if }\protocolsfull \text{ and } k\in I
  $
  \caption{Types for constants}
  \label{fig:type-constants}
\end{figure}

% \\
% \ESelect{C_i} &= \TBForall {(\overline{\alpha\colon\kind})(\tbind\beta\kinds}
% {\oplus(\TPApp P{\overline{\alpha}}). \beta \to  {\overline{!T_i^\pm}. \beta}}
% &
% \offerk &= \TBForall {(\overline{\alpha}\colon K)(\tbind\beta\kinds} {\&
% (\TPApp P {\overline{\alpha}}).\beta \to 
% \TVariant C {{\overline{?T^\pm}}.\beta}}
% \\ 
% & \text{if } \protocols

%%% Local Variables:
%%% mode: latex
%%% TeX-master: "main"
%%% End:

\begin{figure}[t!]
\declrel{Type synthesis}{$\expSynth\Gamma eT\Gamma$}
\begin{mathpar}
  \axiom{\rulenameexpConst}{\expSynth{\Gamma}{c}{\typeof{\termc c}}{\Gamma}}
  
  \infrule{\rulenameexpVar}{\typeAgainst{T}{\kindt}}
  {\expSynth{\Gamma,\ebind xT}{x}{\typec T}{\Gamma}}

  \infrule{\rulenameexpUVar}{\typeAgainst{T}{\kindt}}
  {\expSynth{\Gamma,\eubind xT}{x}{\typec T}{\Gamma,\eubind xT}}

  \infrule{\rulenameexpLetUnit}
  {\expAgainst{\Gamma_1}{e_1}{\TUnit}{\Gamma_2} \\\\
    \expSynth{\Gamma_2}{e_2}{\typec T}{\Gamma_3}
  }
  {\expSynth{\Gamma_1}{\letu {e_1}{e_2}}{\typec T}{\Gamma_3}}
  
  \infrule{\rulenameexpAbs}
  {\typeAgainst{T}{\kindt} \\ \Normal[+]T = \typec V \\\\
    \expSynth{\Gamma_1, \ebind x{V}}{e}{U}{\Gamma_2} \\ \termc{x}\not\in\Gamma_2}
  {\expSynth{\Gamma_1}{\abse xTe}{\TFun[]{V}{U}}{\Gamma_2}}

  \infrule{\rulenameexpApp}
  {\expSynth{\Gamma_1}{e_1}{\TFun[]{T}{U}}{\Gamma_2}\\
    \expAgainst{\Gamma_2}{e_2}{T}{\Gamma_3}} 
  {\expSynth{\Gamma_1}{\appe {e_1}{e_2}}{\typec U}{\Gamma_3}}
  \quad
  \infrule{\rulenameexpPair}{\expSynth{\Gamma_1}{e_1}{T}{\Gamma_2} \\
  \expSynth{\Gamma_2}{e_2}{U}{\Gamma_3}} 
  {\expSynth{\Gamma_1}{\paire {e_1}{e_2}}{\TPair {\typec T} {U}}{\Gamma_3}}

  \infrule{\rulenameexpRec}
  {\typeAgainst{T}{\kindt} \\
    \Normal[+]{\TFun[]TU}=\typec V \\
    \expAgainst{\Gamma, \eubind x{V}}{v}{{V}}{\Gamma, \eubind x{V}}}
  {\expSynth{\Gamma}{\rece x{{\TFun[]TU}}v}{V}{\Gamma}}

  \infrule{\rulenameexpTAbs}
  {\expSynth[\Delta, \typec{\alpha\colon \kind}]{\Gamma_1}{v}{\typec T}{\Gamma_2}} 
  {\expSynth{\Gamma_1}{\tabse \alpha\kind v}{\TForall \alpha\kind T}{\Gamma_2}}

  \infrule{\rulenameexpTApp}
  {\expSynth{\Gamma_1}{e}{\TForall \alpha\kind U}{\Gamma_2} \\ \typeAgainst{T}{\kind}} 
  {\expSynth{\Gamma_1}{\tappe e T}{\Normal[+]{\typec U\tsubs{T}{\alpha}}}{\Gamma_2}}

  % \infrule{\rulenameexpPair}{\expSynth{\Gamma_1}{e_1}{T}{\Gamma_2} \\
  %   \expSynth{\Gamma_2}{e_2}{U}{\Gamma_3}} 
  % {\expSynth{\Gamma_1}{\paire {e_1}{e_2}}{\TPair {\typec T} {U}}{\Gamma_3}}

  \infrule{\rulenameexpLet}
  {\expSynth{\Gamma_1}{e_1}{\TPair TU}{\Gamma_2}\\
    \expSynth{\Gamma_2, \ebind x{\Normal[+]T}, \ebind y{\Normal[+]U}}{e_2}{V}{\Gamma_3}\\
    \termc{x},\termc{y}\not\in\Gamma_3} 
  {\expSynth{\Gamma_1}{\lete xy{e_1}{e_2}}{\typec V}{\Gamma_3}}

  \infrule{\rulenameexpMatch}{
    %\Normal[+]{S}={\TIn{(\TPApp \proto {\overline U})}{S'}} \\
    \protocolsfull \\
    \expSynth {\Gamma_1} e {\TIn{(\TPApp \proto {\overline U})}{S}} {\Gamma_2} \\
    \forall i\in I: \ \expSynth {\Gamma_2,\ebind{x_i}{\matchin}} {e_i} {V_i} {\Gamma_i} \text{  with  }
    \termc{x_i}\not\in\Gamma_i \\
    % MW
    \hl{\typec{V} = \bigsqcup_{i \in I} \typec{V_i}} \\
    % \isEquiv{\Gamma_3}{\Gamma_i}
    % MW
    % \hl{{\Gamma_3} = \bigsqcup_{i \in I} {\Gamma_i}}
  }{
    \expSynth {\Gamma_1} {\matchwithe eCxeiI} {V} {\Gamma_3}
  }
\end{mathpar}
\declrel{Type against}{$\expAgainst\Gamma eT\Gamma$}
\begin{mathpar}
  \infrule{\rulenameexpCheck}
  % MW
  {\expSynth{\Gamma_1} e U {\Gamma_2} \\ \hl{\typec U \subt \typec T}}
  {\expAgainst{\Gamma_1} e T {\Gamma_2}}  
\end{mathpar}
\caption{Algorithmic typing rules for expressions}
\label{fig:alg-typing}
\end{figure}

  % \infrule{\rulenameexpSelect}{
  %   \protocolsfull
  % }{
  %   \expSynth \Gamma {\ESelect{C_i}}
  %   {\TForall{\overline{\alpha}}{\overline\kindp}
  %     {
  %       {\TForall \sigma \kinds {\TFun[]
  %         {\TOut{(\TPApp \proto {\overline\alpha})}\sigma}
  %         {\polarizedparam{!}{\overline{T_i}}\sigma}
  %       }}}}
  %   \Gamma
  % }

  % \infrule{\rulenameexpNew}
  % {\typeAgainst{T}{\kinds}} 
  % {\expSynth{\Gamma}{\newe T}{\TPair {\typec T} {\TDual T}}{\Gamma}}
  
  % \infrule{\rulenameexpUnVar}
  % {\typeSynth{T}{\kindc{v}^\un}}
  % {\expSynth{\Gamma,\termc{x}\colon\typec{T}}{x}{T}{\Gamma, \termc{x}\colon\typec{T}}}
  % 
  % \infrule{\rulenameexpUnAbs}
  % {\typeAgainst{T}{\kindt}
  %   \\ \expSynth{\Gamma_1, \termc{x}\colon\typec{T}}{e}{U}{\Gamma_2}
  %   \\ \Gamma_1 = \Gamma_2 \div \termc{x}}
  % {\expSynth{\Gamma_1}{\abse[\un]xTe}{\TFun[\un]{T}{U}}{\Gamma_1}}
  % 
% \pt{Potentially obsolete rules}
% \begin{mathpar}
%     \infrule{\rulenameexpCtor}
%   {\protocolsdata }
%   {\expSynth{\Gamma}{\constr_i}{\TForall{\overline{\alpha}}{\overline{\kindt}}
%       {\TFun[]{T_{i1}}{\dots \TFun[]{T_{ik}}{\typec{D\
%               \overline{\alpha}}}}}}{\Gamma}}

%   \infrule{\rulenameexpCase1}
%   {
%     \expSynth\Gamma e {\typec{D\ \overline{U}}} {\Gamma_0} \\
%     \protocolsdata \\
%     \expSynth{\Gamma_0, \termc{\overline{x_i}}\colon
%       \typec{\overline{T_i\tsubs{\overline{U}}{\overline{\beta}}}}}{e_i}{\typec V_i}{\Gamma_i}
%     \\
%     \typec V =_\alpha \Normal[+]{V_i}
%     \\
%     \Gamma' =_\alpha \Normal[+]{\Gamma_i \div \overline{x_i}}
%   }
%   {\expSynth{\Gamma}{\casee e {\{\constr_i\ \overline{x_i} \rightarrow
%         e_i\}}}{\typec V}{\Gamma'}}
% \end{mathpar}
%   \infrule{\rulenameexpCase1}{
%     \expSynth \Gamma e {S_0} {\Gamma_0} \\
%     % \typec{S_0} \neq \typec{\proto\ \overline{U}} \\
%     \Normal[+]{S_0} = \typec{?(\TPApp \proto {\overline{U}}).S} \\
%     \protocolsdata \\
%     \expSynth {\Gamma_0,\termc{x_i}\colon {
%     \polarizedparams{?}{\overline{T_i\tsubs{\overline U}{\overline\alpha}}}{S}
%     %?(\typec{\overline{T_i^\pm\tsubs{\overline G}{\overline\alpha}}.S}
%     })
%   } {e_i} {V_i} {\Gamma_i} \\
%   \Normal[+]{V_1} =_\alpha \Normal[+]{V_i} \\
%   \Normal[+]{\Gamma_1} =_\alpha \Normal[+]{\Gamma_i}
%   }{
%     \expSynth \Gamma {\casee e {C_i\,x_i \to e_i, \dots}} {\Normal[+]V} {\Normal[+]{\Gamma'}}
%   }

%   \infrule{\rulenameexpSelect}{
%     \protocolsdata
%   }{
%     \expSynth \Gamma {\ESelect{C_i}}
%     {\TForall{\overline{\alpha}}{\overline\kind}
%       {
%         {\TForall \sigma \kinds {\TFun[]
%           {\typec{!}(\TPApp \proto {\overline\alpha}).\sigma}
%           {\polarizedparam{!}{\overline{O_i}}.\sigma}
% %          {\overline{\blk{!}T_i^\pm}}.S
%         }}}}
%     \Gamma
%   }


%%% Local Variables:
%%% mode: latex
%%% TeX-master: "main"
%%% End:


Expression typing is again mostly standard with respect to (linear) functional
or session type languages. The definitions for typing assume a set of
protocol declarations as an implicit argument. The typings of
$\appe\selectk{C}$ (\cref{fig:type-constants}) and $\matchk$
(\cref{fig:alg-typing}, \rulenameexpMatch) access these protocol declarations. 
The $\typeofop$ operator in~\cref{fig:type-constants} returns the
types for constants in normal form. The typing rules for expressions are defined
in~\cref{fig:alg-typing}. 
%
The type system is algorithmic, implemented with two mutually recursive
judgments: one to synthetize a type of an expression, the other to check an
expression against a type. The typing rules rely on type variable contexts
$\Delta$ introduced in \cref{sec:types}, on term variable contexts $\Gamma$
assigning types $\typec T$ to term variables $\termc x$. There are two sorts of
entries in term variable contexts: linear entries, noted $\ebind xT$, and
unrestricted entries, noted $\eubind xT$. The latter are used for recursive
functions only (\cf rules \rulenameexpUVar and \rulenameexpRec) and are inspired
by work on quantitative type theory
\cite{DBLP:conf/lics/Atkey18,DBLP:conf/birthday/McBride16,DBLP:journals/pacmpl/BernardyBNJS18}.

The type
system maintains the invariant that entries in term variable contexts are always normalized:
the rules that write into contexts, namely
\rulenameexpAbs, \rulenameexpRec and \rulenameexpLet, normalize the
new entry. 
Judgments are of the form $\expSynth{\Gamma_1} eT{\Gamma_2}$ or
$\expAgainst{\Gamma_1} eT{\Gamma_2}$, where $\Gamma_2$ represents the
part of context $\Gamma_1$ \cite{walker:substructural-type-systems}
that is not consumed by $\termc e$ and $\typec T$ is in normal form.
% or the unused part of a session type for a given endpoint
% \cite{DBLP:journals/iandc/Vasconcelos12,DBLP:conf/forte/ZalakainD21}.
The axioms
(\rulenameexpConst and \rulenameexpVar) mostly copy the input context to the
output. Rule \rulenameexpApp is a good example of context threading: 
to type expression $\appe {e_1}{e_2}$, expression $\termc{e_1}$ is given the
input context $\Gamma_1$, produces $\Gamma_2$ which is given to $\termc{e_2}$,
which turn produces $\Gamma_3$ that constitutes the outgoing context of
application $\appe {e_1}{e_2}$.
%
The rules that add linear entries to the context (\rulenameexpAbs, \rulenameexpLet
and \rulenameexpMatch) all check that the new variables are not present
in the outgoing context, thus ensuring that the corresponding values are fully
consumed. 
% 
Rule \rulenameexpTApp explicitly normalizes the output type for normalization is
not preserved by substitution; all rules that introduce new entries in the
context (\rulenameexpAbs, \rulenameexpRec, \rulenameexpLet, and
\rulenameexpMatch) make sure that types are normalized, thus ensuring that contexts contain only
normalized types.
% other rules
% construct types in normal form, either by assumption or by induction on a
% premise (see \cref{sec:appendix-exps-procs}).

% ======= I don't think we need this
% Rule \rulenameexpTApp normalizes the output type, all other rules
% construct types in normal form, either by assumption or by induction
% on a premise (see \cref{sec:appendix-exps-procs}, \cref{lem:synth-against-nf}).
% >>>>>>> e9014ece406b68d500fbdef7db6faf10a460f9ae

Rules \rulenameexpRec and \rulenameexpUVar govern the typing of
recursive functions. Rule \rulenameexpRec binds the recursive function $x$
using an unrestricted binding $\eubind x\_$ so that any number of
recursive calls is possible in the body via \rulenameexpUVar. Unrestricted
bindings provide for termination of recursive functions. Such functions cannot use linear variables from the
environment, which we enforce by requiring the incoming context
$\Gamma$ to be equal to the outgoing context.

The truly new rule is \rulenameexpMatch, which indicates that all
rules are parametric with respect to an implicit set of protocol
declarations: there is a declaration for each protocol $\proto$ which
has a kinding in $\Delta$. The rule starts by synthesizing the 
input type of the expression to be matched;
the type is 
$\TIn{(\TPApp \proto {\overline U})}{S}$ for some
protocol constructor $\proto$. At this point, the declaration of protocol
$\proto$ is looked up in the implicit set of declarations. The
$\matchk$ branch for each choice of selector $\termc{C_i}$ is then typed in turn. The
parameters $\typec{\overline \alpha}$ of the protocol are instantiated with
the arguments $\typec{\overline U}$ to produce a sequence of types
$\typec{\overline{T_i}\tsubs{\overline U}{\overline \alpha}}$. This sequence of
types is converted to a normalized session type as described in \cref{sec:types},
which is assigned to variable $\termc{x_i}$. Now, each branch synthesizes a type
$\typec{V_i}$ and context $\Gamma_i$, for $i\in I$. 
\hl{[\text{MW}]:} In all branches,
the contexts must coincide (up to $\alpha$
equivalence) with the output context of the $\matchk$. 
The output type $\typec V$ is the least upper bound of the branch types. 

% Similarly, the context $\Gamma_3$ is computed via a pointwise join: $\forall x \in \Gamma_3$, 
% $\Gamma_3(x) = \bigsqcup_{i \in I} \Gamma_i(x)$.
% OLD:
% In all branches,
% the output types and the contexts must coincide (up to $\alpha$
% equivalence) with the output type 
% and context of the $\matchk$, i.e., $\typec V$ and $\Gamma_3$. In practice,
% the implementation selects one branch, say $k\in I$, and compares all
% other output types and contexts with $\typec{V_k}$ and $\Gamma_k$.
In this rule, and also in the type for $\selectk$ in \cref{fig:type-constants}, we extend the
directional operators and  materialization (from \cref{fig:type-polarized-operator}) 
to sequences
of parameters by mapping:
$\polarizedparam{\pm}{\overline{T}}{} =
\overline{\polarizedparam{\pm}{T}{}}$ as well as
$\tosession {\varepsilon}S = \typec S$ and $\tosession {T\overline T}S =
\tosession{T}{\tosession {\overline T}S}$.

\hl{[\text{MW}]:} Rule \rulenameexpCheck relies on the previous judgement to check expression $\termc e$ against 
type $\typec T$, up to subtyping.
Both types can be assumed to have normal form, because synthesis outputs $\typec U$ in normal form and
checking must be invoked with $\typec T$ in normal form.

% OLD:
% Rule \rulenameexpCheck relies on the previous judgement to check expression $\termc e$ against 
% type $\typec T$, up to $\alpha$-equivalence. This simple comparison is
% sufficient because synthesis outputs $\typec U$ in normal form and
% checking must be invoked with $\typec T$ in normal form.

% \begin{lemma}\
%      \label{lem:synth-against-nf}
%      \begin{enumerate}
%           \item\label{it:synth-nf} If $\expSynth{\Gamma_1} eT{\Gamma_2}$ and $\isCtx[\Delta]{\Gamma_1}{\kindt}$, 
%           then $\isCtx[\Delta]{\Gamma_2}{\kindt}$ and $\typec T$ is in normal form.
%           \item\label{it:against-nf}  If $\expAgainst{\Gamma_1} eT{\Gamma_2}$ and $\isCtx[\Delta]{\Gamma_1}{\kindt}$
%           and $\typec T$ is in normal form, then $\isCtx[\Delta]{\Gamma_2}{\kindt}$.
%      \end{enumerate}
%      \begin{proof}\
%      \begin{enumerate}
%           \item By rule induction on the first hypothesis.
%           \item Using \cref{it:synth-nf} and observing that type $\typec U$ such that 
%           $\typec U =_\alpha \typec T$ is also in normal form.
%           \qedhere
%      \end{enumerate}         
%      \end{proof}
% \end{lemma}
%
%
% \paragraph*{Operational semantics}

The operational semantics for expressions and processes is defined via a 
transition relation labelled by the below actions.
%
\begin{align*}
  \sigma &\grmeq \receivel xv
	   \grmor \sendl xv
       \grmor \receivel xC
	   \grmor \sendl xC
       \grmor \closel x
	   \grmor \openl x
       && \text{labels for session operations}
  \\
  \lambda &\grmeq \sigma
       \grmor \beta
       \grmor \forkl v
	   \grmor \newl xxT
       && \text{labels for expressions}
  \\
  \pi &\grmeq \sigma 
  \grmor \tau
  \grmor \scopel xxxx
  \grmor (\parl \pi \pi)
       && \text{labels for processes}
\end{align*}
%
The $\sigma$ labels capture the six operations on sessions: $\receivek$,
$\sendk$, $\matchk$, $\selectk$, $\waitk$ and $\terminatek$.
%
Labels $\lambda$ are for expressions: $\beta$ is for internal actions and the
remaining two for $\forkk$ and $\newk$.
%
Labels $\pi$ are for processes: $\tau$ is for silent actions (the counterpart of
$\beta$ for expressions), $\scopel abxa$ or $\scopel abxb$ is for opening the
scope of a $\termc{(\nu ab)}$ binding, and $\pi\PAR\pi$ for parallel
composition.

\begin{figure}[t!]
  \declrel{Labelled transition system for expressions}{$\trans e \lambda e$}
  \begin{mathpar}
    \inferrule[Act-App]{}{\trans{\appe{(\abse{x}{T}{e})}{v}}{\beta}{\termc{e}\esubs{v}{x}}}
    
	\inferrule[Act-TApp]{}{\trans{\tappe{(\tabse{\alpha}{\kind}{v})}{T}}{\beta}{\termc{v}\tsubs{T}{\alpha}}}

	\inferrule[Act-Let]{}{\trans{\lete{x}{y}{{\paire uv}}{e}}{\beta}{\termc{e}\esubs{u}{x}\esubs{v}{y}}}

	\inferrule[Act-Let*]{}{\trans{\letu{\unitk}{e}}{\beta}{\termc{e}}}

    \inferrule[Act-Rec]{}{\trans{\appe{({\rece xTv})}
        u}{\beta}{\appe{(\termc{v}\esubs{\rece xTv}{x})}u}}

    \inferrule[Act-Fork]{}{\trans{\appe{\forkk}{v}}{\forkl{v}} \unitk}
	
	\inferrule[Act-New]{}{\trans{\tappe \newk T}{\newl xyT}{\paire xy}}

	\inferrule[Act-Rcv]{}{\trans{\appe{\tappe{\tappe \receivek T}{U}}{x}}{\receivel xv}{\paire{v}{x}}}
	
	\inferrule[Act-Send]{}{\trans{\appe{\appe{\tappe{\tappe \sendk T}{U}}{v}}{x}}{\sendl xv}{x}}
	
	\inferrule[Act-Match]{}{\trans{\matchwithe xCyeiI}{\receivel
        x{C_k}}{e_k\esubs x{y_k}}}

	\inferrule[Act-Sel]{}{\trans{\appe{\tappe{\selecte{C}}{\overline T}}{x}}{\sendl x{C}}{x}}
	
	\inferrule[Act-Wait]{}{\trans{\waite{x}}{\closel x}\unitk}

	\inferrule[Act-Term]{}{\trans{\closee{x}}{\openl x}{\unitk}}
  \end{mathpar}
  \caption{Labelled transition system for expressions (selected rules). The full version, including
  call-by-value structural rules, is presented in \cref{sec:appendix-exps-procs}, \cref{fig:lts-exps}}
  \label{fig:lts-exps-abbrev}
\end{figure}

%%% Local Variables:
%%% mode: latex
%%% TeX-master: "main"
%%% End:


The axioms of the labelled transition system for
expressions is defined by the rules in \cref{fig:lts-exps-abbrev}
(\cref{sec:appendix-exps-procs}, \cref{fig:lts-exps}, contains the whole set).
%
We group under label $\beta$ the common reduction rules of the polymorphic
lambda calculus. Rules Act-Fork and Act-New record in the label the value forked
and the names and types of the channel ends created. These two transitions will
then be handled at the process level. Then we have the six transitions for the
common session operations.
% These are now in appendix
% Lastly we have the call-by-value structural rules.
% For example, rule Act-AppL reduces the left hand side of an application, whereas
% Act-AppR reduces the right hand side when the left is a value.

\begin{figure}[t!]
  \declrel{Labelled transition system for process configurations}
    {$\trans K\pi{K'}$}
  \begin{mathpar}
    \infrule{Act-Beta}
    {\trans{e}{\beta}{e'} }
    {\trans{(\{e\}\uplus E,L)}\tau
      {(\{e'\}\uplus E,L)} }

    \infrule{Act-Fork}
    {\trans{e}{\forkl v}{e'} }
    {\trans{(\{e\}\uplus E,L)}\tau
      {(\{e',\appe v\unitk\}\uplus E,L)} }

    \infrule{Act-New}
    {\trans{e}{\newl xyS}{e'} \and x\ne y \\
      x,y\notin(\{e\}\uplus E,L) }
    {\trans{(\{e\}\uplus E,L)}{\newl xyS}
      {(\{e'\}\uplus E,L\cup\{\{x,y\}\})} }

    \infrule{Act-Msg}
    {\{x,y\}\in L \\
      \trans{e_1}{\receivel xv}{e_1'} \and
      \trans{e_2}{\sendl yv}{e_2'} }
    {\trans{(\{e_1,e_2\}\uplus E,L)}\tau
      {(\{e_1',e_2'\}\uplus E,L)} }

    \infrule{Act-Bra}
    {\{x,y\}\in L \\
      \trans{e_1}{\receivel x{C}}{e_1'} \and
      \trans{e_2}{\sendl y{C}}{e_2'} }
    {\trans{(\{e_1,e_2\}\uplus E,L)}\tau
      {(\{e_1',e_2'\}\uplus E,L)} }

    \infrule{Act-Wait}
    {\{x,y\}\in L \\
      \trans{e_1}{\closel x}{e_1'} \and
      \trans{e_2}{\closel y}{e_2'} }
    {\trans{(\{e_1,e_2\}\uplus E,L)}\tau
      {(\{e_1',e_2'\}\uplus E,L\setminus\{\{x,y\}\})} }
  \end{mathpar}
  \caption{Labelled transition system for process configurations}
  \label{fig:lts-procs}
\end{figure}

%%% Local Variables:
%%% mode: latex
%%% TeX-master: "main"
%%% End:


The labelled transition system for processes is defined by the rules in
\cref{fig:lts-procs}. The rules are adapted from
\citet{DBLP:conf/concur/0001KDLM21,DBLP:journals/corr/abs-2105-08996} which in
turn are inspired by \citet{DBLP:journals/corr/abs-2106-11818}. The first four
rules convert expression transitions into process transitions: the session
transitions $\sigma$ are passed directly from expressions to processes (rule
Act-Session). The next three rules all yield silent ($\tau$) transitions:
Act-Fork launchs a new thread and Act-New creates a new channel. Rules Act-JoinL
and Act-JoinR account for the commutative nature of parallel composition. Rules
Act-ParL, Act-ParR and Act-Res allow transitions underneath parallel composition
and scope restriction. The rules for opening and closing a scope (left and right
versions) are adapted to the double binder scheme from the labelled transitions
systems for the $\pi$-calculus~\cite{DBLP:journals/iandc/MilnerPW92b}. These
rules are required for we work with free output, unlike the above cited
works~\cite{DBLP:conf/concur/0001KDLM21,DBLP:journals/corr/abs-2105-08996,DBLP:journals/corr/abs-2106-11818}
that work with bound output.


Reduction underneath a prefix (rule Act-Res) can only occur if the channels ends
$\termc x$ and $\termc y$ (two bound variables) do not capture the free
variables in the label. For this purpose, we define the free variables of a
process label $\pi$ as follows.
%
\begin{align*}
  % sigma labels
  \FVL {\receivel xv} = \FVL {\sendl xv} = \{x\} \cup \FVE v
  \qquad
  \FVL {\receivel xC} = \FVL {\sendl xC} = \FVL {\closel x} = \FVL {\openl x} = \{x\}
  \\
  \FVL \tau = \emptyset
  \qquad
  \FVL {\scopel wxyz} = \{y,z\} \setminus \{w,x\}
  \qquad
  \FVL {\parl \pi {\pi'}} = \FVL \pi \cup \FVL {\pi'}
\end{align*}


%%% Local Variables:
%%% mode: latex
%%% TeX-master: "main"
%%% End:
